% witness: 2201.01339, 2201.11150 (in IEEEeqnarray cells), 1410.7694 (`\IEEEQEDhere`); 1903.06134 (`\IEEEQEDhereeqn` in an
%          `equation*` in the outer proof of a nested pair, :3279-3281), 2201.03502; 2001.04812 (amsthm `\qedhere` in
%          an IEEEproof)
% oracle:  pdflatex 0 errors (IEEEtran.cls:5549-5556: `\IEEEQEDhere` = `\global\@IEEEQEDshowfalse\IEEEQED`,
%          `\IEEEQEDhereeqn` the same with the QED as the equation number, `\endIEEEproof` typesets `\IEEEQED` while
%          the switch is true); the inner proof's QED where `\IEEEQEDhere` stands and none at the outer's end (the
%          false is global), amsthm's box kept after an `\IEEEQEDhere` in it, the numbered display's one tag "■(1)",
%          twice "■■(2)", IEEEproof's end QED kept after amsthm's `\qedhere`, a document's `\IEEEQED` ("X") at the
%          end, the IEEEeqnarray `x` cell's ■ (2201.11150 :1366, :2285; 1903.06134 treepin.tex:2217-2226), an
%          eqnarray `\mbox{\IEEEQED}`
% engines: rust 64h4 = 5 errors (`\IEEEQEDhereeqn`, `\IEEEQEDoff`, `\@IEEEQEDshowfalse`, `\if@IEEEQEDshow` undefined;
%          the `{rCl+x*}` row "Extra alignment tab"), the
%          nested pairs mis-nested (outer proofs never closed), the amsthm proof's own QED taken by `\IEEEQEDhere`
%          (Perl IEEEtran.cls.ltxml:193-196 pops the shared QED@stack), the eqnarray `\mbox{\IEEEQED}` cell empty (a
%          size-less whatsit makes the cell "empty", Perl Alignment.pm:458-470); Perl the same 5 errors
% expect:  0 errors, 0 warnings; "Inner. ∎", "OuterEnd." without ∎; "Ams. ∎∎"; the `$$` and `equation*` displays
%          tagged ∎, no end QED; the numbered `equation`'s one `ltx:tags`, its displayed tag ∎ then "(1)" (refnum 1,
%          `\eqref{e}` "(1)"), twice ∎∎ then "(2)"; the align proof ends with ∎; the last proof ends "X"; the `x` cell
%          and the `\mbox` cell hold ∎; the class's own switch (IEEEtran.cls:5543 `\newif\if@IEEEQEDshow`): A1 (global
%          false), A2 (local false), A4 (`\IEEEQEDoff`) and A5's outer end without ∎, A3 and A6 (the local false undone,
%          each proof resets it) and A5's inner with it, "SHOWFALSE" after a global false
% status:  RED (64h4) -> GREEN (64ra)
% preload: ar5iv.sty
\documentclass{IEEEtran}
\usepackage{amsmath}
\usepackage{amsthm}
\begin{document}
\begin{IEEEproof}
Outer.
\begin{IEEEproof}
Inner. \IEEEQEDhere
\end{IEEEproof}
OuterEnd.
\end{IEEEproof}
\begin{proof}
Ams. \IEEEQEDhere
\end{proof}
\begin{IEEEproof}
Hence
$$ x = y \IEEEQEDhereeqn $$
\end{IEEEproof}
\begin{IEEEproof}
Outer.
\begin{IEEEproof}
Inner.
\end{IEEEproof}
Thus
\begin{equation*}
a = b. \hfill\IEEEQEDhereeqn
\end{equation*}
\end{IEEEproof}
\begin{IEEEproof}
Numbered
\begin{equation}\label{e}
c = d \IEEEQEDhereeqn
\end{equation}
\end{IEEEproof}
\begin{IEEEproof}
Twice
\begin{equation}
g = h \IEEEQEDhereeqn \IEEEQEDhereeqn
\end{equation}
\end{IEEEproof}
\begin{IEEEproof}
Cell
\begin{IEEEeqnarray*}{rCl+x*}
a &=& b. & \hfill\IEEEQEDhere
\end{IEEEeqnarray*}
\end{IEEEproof}
\begin{eqnarray*}
a &=& \mbox{\IEEEQED}
\end{eqnarray*}
See \eqref{e}.
\begin{IEEEproof}
Then
\begin{align*}
e &= f. \qedhere
\end{align*}
\end{IEEEproof}
{\renewcommand\IEEEQED{X}
\begin{IEEEproof}
Own.
\end{IEEEproof}}
\makeatletter
\begin{IEEEproof} A1. \global\@IEEEQEDshowfalse
\end{IEEEproof}
\begin{IEEEproof} A2 local. \@IEEEQEDshowfalse
\end{IEEEproof}
\begin{IEEEproof} A3 plain.
\end{IEEEproof}
\begin{IEEEproof} A4 off. \IEEEQEDoff
\end{IEEEproof}
\begin{IEEEproof} A5 outer
\begin{IEEEproof} A5 inner \hfill\IEEEQEDhere
\end{IEEEproof}
A5 outer tail.
\end{IEEEproof}
After: \if@IEEEQEDshow SHOWTRUE\else SHOWFALSE\fi.
\begin{IEEEproof} A6 reset.
\end{IEEEproof}
\makeatother
\end{document}
